\songtitle{Marchan las hormigas}

% -> stories.tex (More); original.tex (Other Narratives); generos.tex (The Orchestral & Anthem Wing)
\tag*{2026}\tag{march}\tag{original-lyrics}\tag{spanish-lyrics}
\begin{notes}
Written by Carlos Thompson — a militar march sung by ants, each one certain of her value and her destiny within the colony. Retrieved from Dvigitt's Suno page (V6-MINI, 26 September 2026).
\end{notes}
\begin{style}
Militar march at a steady, buoyant pace, female choir in broad layered harmonies with crisp unison calls, marching percussion, rounded bass, bright woodwinds and rhythmic handclaps; full, balanced choral production.
\end{style}
\begin{lyrics}
\songpart{Intro}
¡Listos!\\
¡Con compás!\\
¡Maaarrrch!

\songpart{Verse 1}
Las hormigas marchan, marchan, marchan\\
cada una sabe, sabe, sabe\\
cual es su valor, lor-lor, lor-lor\\
cual es su destiiiiiino.

\songpart{Chorus}
Marchan las hormigas (marchan las hormigas),\\
tienen su misión (tienen su misión),\\
vivir por la colonia (vivir por la colonia),\\
darlo todo bajo el sol.

\songpart{Verse 2}
Todas sincopadas, padas, padas\\
a cumplir la misión, sióon, sióon\\
de buscar la fuente, fuente, fuente\\
para cultivaaaaaaaaaar

\songpart{Chorus}
Marchan las hormigas (marchan las hormigas),\\
tienen su misión (tienen su misión),\\
vivir por la colonia (vivir por la colonia),\\
darlo todo bajo el sol.

Marchan las hormigas (marchan las hormigas),\\
tienen su misión (tienen su misión),\\
vivir por la colonia (vivir por la colonia),\\
darlo todo bajo el sol.
\end{lyrics}

